\documentclass[11pt]{article}
\usepackage[T1]{fontenc}
\usepackage{textcomp}
\usepackage[margin=0.75in]{geometry}
\usepackage{amsmath,amssymb,amsthm}
\usepackage{array,booktabs}
\usepackage{hyperref}
\setlength{\parindent}{0pt}
\newtheorem{theorem}{Theorem}
\theoremstyle{definition}
\newtheorem{definition}{Definition}
\title{Flow Proof Artifact: Comb}
\author{Generated by \texttt{flow doc proof}}
\date{Flow Proof Book}
\begin{document}
\maketitle
Basic combinatorial identities on binomial coefficients.\\[0.5em]
\textbf{Source.} graham-knuth-patashnik — *Concrete Mathematics*\\[0.5em]
\bigskip
\noindent{\large \textbf{Definition 1.} \textit{Choosing zero items always gives one way} \hfill Comb \cdot choose \cdot choosing none gives one}\par
\smallskip\noindent\textit{Coordinate: Comb \cdot choose \cdot choosing none gives one}\par
\smallskip\noindent\textit{Source: graham-knuth-patashnik}\par
\medskip
\noindent\textbf{Goal.} Choosing zero items always gives one way\par
\noindent$\displaystyle \forall n \in \mathbb{N}\quad choose(n, 0) = 1$\par
\medskip
\noindent
\begin{tabular}{@{} >{\raggedright\arraybackslash}p{0.06\textwidth} >{\raggedright\arraybackslash}p{0.47\textwidth} >{\raggedleft\arraybackslash}p{0.41\textwidth} @{}}
\textbf{\#} & \textbf{Proof} & \textbf{Mathematics} \\
\midrule
\textbf{1.} & We stipulate choosing none gives one for choose on Comb — this is a definition, not a derived fact. & \\[0.45em]
\textbf{2.} & This follows directly from the definition: choose(n, 0) equals 1. Hence proven. & $\displaystyle choose(n, 0) = 1$ \\[0.45em]
\bottomrule
\end{tabular}
\label{thm:Comb-choose-choosing_none_gives_one}
\par\medskip\hrule
\bigskip
\noindent{\large \textbf{Definition 2.} \textit{Pascal recurrence for binomial coefficients} \hfill Comb \cdot choose \cdot Pascal recurrence holds}\par
\smallskip\noindent\textit{Coordinate: Comb \cdot choose \cdot Pascal recurrence holds}\par
\smallskip\noindent\textit{Source: graham-knuth-patashnik}\par
\medskip
\noindent\textbf{Goal.} Pascal recurrence for binomial coefficients\par
\noindent$\displaystyle \forall n \in \mathbb{N} \forall k \in \mathbb{N}\quad choose(n, k) = choose(pred(n), j) + choose(pred(n), k)$\par
\medskip
\noindent
\begin{tabular}{@{} >{\raggedright\arraybackslash}p{0.06\textwidth} >{\raggedright\arraybackslash}p{0.47\textwidth} >{\raggedleft\arraybackslash}p{0.41\textwidth} @{}}
\textbf{\#} & \textbf{Proof} & \textbf{Mathematics} \\
\midrule
\textbf{1.} & We stipulate pascal recurrence holds for choose on Comb — this is a definition, not a derived fact. & \\[0.45em]
\textbf{2.} & We split into exhaustive cases — the claim must hold in each one. & \\[0.45em]
\textbf{3.} & Case 1 (see step 2): suppose k is zero. & \\[0.45em]
\textbf{4.} & From step 3, this implies choose(n, k) equals 1 in this case. & $\displaystyle choose(n, k) = 1$ \\[0.45em]
\textbf{5.} & Case 2 (see step 2): neither disjunct holds. & \\[0.45em]
\textbf{6.} & Let j denote the predecessor of k. & \\[0.45em]
\textbf{7.} & From step 5 and step 6, this implies choose(n, k) equals choose(pred(n), j) plus choose(pred(n), k). Together with the other cases (step 3 and step 5), the goal is discharged. Hence proven. & $\displaystyle choose(n, k) = choose(pred(n), j) + choose(pred(n), k)$ \\[0.45em]
\bottomrule
\end{tabular}
\label{thm:Comb-choose-pascal_recurrence_holds}
\par\medskip\hrule
\end{document}
